\begin{theorem}[CertificateCompiler bridge schema handoff]
\label{thm:certificate-compiler-bridge-schema-handoff}
Every public $\CertificateCompilerUp$ packet hands a schema bridge only the finite compiler graph, target landing rows, classifier-transport rows, continuation rows, package provenance, and source, target, and local $\NameCert$ rows. The handoff is the displayed packet surface $S,T,G,E,C,H,P,N$ together with its identity, composite, and triple-composite ledgers; it does not include host compiler semantics, optimizer correctness, decompiler behavior, ambient program equality, or unrecorded target behavior.
\end{theorem}
\begin{proof}
Start with the carrier of \autoref{def:certificate-compiler-carrier}. It exposes source and target certificates, graph edges, landing rows, continuation rows, componentwise transports, provenance, and a local naming row.

Use the classifier and the NameCert obligations. \autoref{def:certificate-compiler-classifier} and \autoref{thm:certificate-compiler-namecert-obligations} identify the bridge-readable packet with the finite graph, landing, classifier transport, continuation, package, and naming rows.

Now apply edge-classifier exactness, the identity ledger, composition stability, and triple-composite associativity. These theorems show that every schema-facing edge, identity, binary composite, and triple composite is replayed through the same finite rows, so the bridge handoff has no host compiler coordinate.
\end{proof}

\begin{theorem}[CertificateCompiler bridge-row non-escape]
\label{thm:certificate-compiler-bridge-rows-nonescape}
The bridge-facing rows of a public $\CertificateCompilerUp$ packet export no optimizer, decompiler, program equality, ambient categorical equality, or unrecorded target behavior. Any schema bridge that consumes the packet must read the source certificate, target certificate, graph edge, landing, route, transport, provenance, and local naming rows exactly as displayed.
\end{theorem}
\begin{proof}
Project the public packet as the finite rows $S,T,G,E,C,H,P,N$. These rows are the only bridge-facing coordinates of the carrier.

Apply \autoref{thm:certificate-compiler-edge-classifier-exactness} and \autoref{thm:certificate-compiler-finite-route-non-escape}. They force edge reads through graph, landing, classifier transport, continuation, provenance, and naming rows and exclude host compiler semantics from the packet.

The identity ledger, composition stability, and triple-composite associativity preserve the same row set under bridge replay. Since no displayed row is an optimizer, decompiler, program-equality witness, ambient categorical equality, or target-behavior oracle, none can be exported to the bridge consumer.
\end{proof}

\begin{theorem}[CertificateCompiler bridge consumer exhaustion]
\label{thm:certificate-compiler-bridge-consumer-exhaustion}
Every schema-bridge consumer of a public $\CertificateCompilerUp$ packet factors through the already displayed source certificate, target certificate, graph edge, landing row, route row, classifier-transport row, provenance row, and local naming row. The consumer can compose, reassociate, and replay certificate-compiler edges only by the identity, binary-composite, and triple-composite ledgers already present in the chapter.
\end{theorem}
\begin{proof}
Begin with the public consumer boundary \autoref{thm:certificate-compiler-public-consumer-boundary}. It already restricts downstream reads to graph landing, classifier transport, continuation, provenance, and local naming over the source and target certificates.

Add the bridge handoff of \autoref{thm:certificate-compiler-bridge-schema-handoff}. The bridge receives exactly the same finite graph and certificate rows, with no host compiler semantics.

Finally apply the identity ledger, composition stability, and triple-composite associativity. These results exhaust the bridge consumer's available operations on the packet: identity replay, binary composition, and triple reassociation all factor through the displayed rows and cannot introduce another compiler surface.
\end{proof}

\begin{theorem}[CertificateCompiler bridge associativity split index]
\label{thm:certificate-compiler-bridge-associativity-split-index}
The bridge handoff of \autoref{thm:certificate-compiler-bridge-schema-handoff}, the bridge-row non-escape boundary of \autoref{thm:certificate-compiler-bridge-rows-nonescape}, and the bridge consumer exhaustion row of \autoref{thm:certificate-compiler-bridge-consumer-exhaustion} all read the same identity, binary-composite, and triple-composite ledger for $\CertificateCompilerUp$ packets. A schema bridge therefore consumes one finite ledger surface rather than three unrelated compiler interfaces.
\end{theorem}
\begin{proof}
First read the identity ledger. \autoref{thm:certificate-compiler-identity-ledger} gives the diagonal compiler packet over the same displayed source, target, graph, landing, classifier transport, continuation, provenance, and naming rows consumed by the bridge handoff.

Next read binary composition. \autoref{thm:certificate-compiler-composition-stability} and \autoref{thm:certificate-compiler-identity-composite-ledger-exhaustion} keep composite replay on the same public row surface and add no host compiler semantics.

Finally read triple composites. \autoref{thm:certificate-compiler-triple-composite-associativity} and the bridge consumer exhaustion theorem use the same triple of graph edges and middle landing matches, while bridge-row non-escape blocks optimizer, decompiler, program-equality, ambient categorical, and unrecorded target-behavior coordinates.
\end{proof}
